\chapter{PWA Reference Implementation}
\label{appendix:pwa}

\lstdefinelanguage{json}{
    basicstyle=\ttfamily\small,
    showstringspaces=false,
    breaklines=true,
    frame=single,
    morestring=[b]",
    stringstyle=\color{blue},
    comment=[l]{//},
    morecomment=[s]{/*}{*/},
    commentstyle=\color{gray}\itshape
}

To operationalize the entire \textit{Bridging Worlds} framework, this appendix documents a production-ready progressive web application (PWA) that captures the guided workflows, quality instrumentation, collaboration patterns, and knowledge assets discussed throughout the book. The implementation is intentionally modular so teams can adopt the components incrementally.

\section{Architecture Overview}
\begin{itemize}
    \item \textbf{Front end}: React~(TypeScript) with a PWA shell (Workbox/Vite PWA). Redux Toolkit or Recoil handles client state, while React Query synchronizes data with the API. Material UI/Tailwind provide theming. Offline caching stores templates, glossary entries, and recent requirements in IndexedDB/Dexie.
    \item \textbf{Back end}: NestJS or FastAPI backed by PostgreSQL (Prisma/SQLAlchemy). Auth is delegated to Auth0/Azure AD. Background workers (BullMQ/Celery) process long-running quality checks and notifications.
    \item \textbf{Service integrations}: Feature-flag providers (LaunchDarkly/Optimizely) for experiment orchestration, ticketing integrations (Jira/Azure DevOps), and analytics pipelines (Segment/Snowflake) for metric ingestion.
\end{itemize}

\section{Guided Requirement Workflows}
\begin{itemize}
    \item \textbf{Dynamic wizards} mirror the templates from Chapters~\ref{ch:data-informed-tickets}, \ref{sec:experiment-design}, and \ref{ch:documentation-patterns}. JSON schema forms render sections (problem statement, hypotheses, metrics, acceptance criteria, risks, stakeholders) with autosave and validation.
    \item \textbf{Template library} stores curated artifacts (feature, exploration, ML-model, experiment, documentation). Users can fork, version, and diff templates. Inline glossary tooltips link to the definitions introduced in Chapter~\ref{sec:mental-models}.
    \item \textbf{Experiment planning} embeds the logic from \texttt{tools/experiment\_planner.py}. Hypothesis prompts, sample-size calculators, tracking matrices, and analysis report shells are exposed directly in the UI.
\end{itemize}

\section{Integrated Quality and Metrics Tracking}
\begin{itemize}
    \item A quality-analysis microservice ports the logic from \texttt{tools/quality\_analyzer.py} so every requirement receives completeness, clarity, consistency, and acceptance-quality scores.
    \item Dashboards visualize requirement quality trends, experiment status, and metric guardrails (Chapters~\ref{ch:quantifying-quality} and \ref{sec:metrics-dashboard-prototypes}). Alerts notify teams when thresholds are breached.
    \item Metric definitions link to live analytics sources, allowing teams to attach baseline/target/actuals per requirement and monitor progress in real time.
\end{itemize}

\section{Collaboration Features}
\begin{itemize}
    \item \textbf{Real-time co-editing}: CRDT libraries (Yjs/Automerge) enable concurrent editing with presence indicators, threaded comments, and change history.
    \item \textbf{Stakeholder alignment tracking}: The facilitation canvases from \texttt{tools/facilitation.py} are embedded as interactive dashboards to capture influence, alignment, and risk owners (Chapter~\ref{sec:facilitation-tools}).
    \item \textbf{Workflow states}: Requirements move through configurable stages (Draft, Review, Aligned, Approved, Experimenting, Locked). Status changes propagate to connected ticketing systems via webhooks.
\end{itemize}

\section{Knowledge Base and Search}
\begin{itemize}
    \item The knowledge base ingests patterns, templates, and case studies from Chapters~\ref{ch:documentation-patterns}--\ref{ch:templates-frameworks}. Content is stored as Markdown/MDX and indexed via MeiliSearch/Elastic for full-text search.
    \item Contextual helpers surface related entries when users edit hypotheses, acceptance criteria, or documentation patterns, ensuring the framework is always one click away.
    \item Users can bookmark, comment, or propose edits, creating a virtuous cycle of curated knowledge.
\end{itemize}

\section{DevOps and Deployment}
\begin{itemize}
    \item Front end deploys on Vercel/Netlify with CI pipelines (GitHub Actions) running lint, unit, and Playwright/Cypress suites.
    \item Back end runs on container platforms (AWS Fargate/Azure App Service) with managed PostgreSQL, S3/Blob storage for attachments, and queue workers for asynchronous jobs.
    \item Security includes role-based access control (PM, Engineer, Data Scientist, Executive), audit logs, encryption at rest, and SOC2-ready logging.
\end{itemize}

\section{Adoption Roadmap}
\begin{enumerate}
    \item \textbf{MVP}: Requirement wizard, quality analyzer integration, knowledge base reader, and lightweight dashboards.
    \item \textbf{Collaboration \& Experimentation}: Real-time co-editing, stakeholder alignment tracking, experiment planning module.
    \item \textbf{Advanced Analytics}: Metric ingestion, alerting, and integrations with external experiment platforms.
    \item \textbf{Ecosystem}: Plugin APIs, AI-assisted prompts, deeper ticketing/documentation sync.
\end{enumerate}

\section{Integration Blueprint}

The PWA integrates with the broader toolchain through concrete APIs.

\subsection{Requirement Quality Metrics to Jira/GitHub}
\paragraph{Jira comments.} A worker posts quality summaries as Atlassian “doc” comments via \texttt{/rest/api/3/issue/\{id\}/comment}:
\begin{lstlisting}[language=json]
POST https://your-domain.atlassian.net/rest/api/3/issue/REQ-123/comment
Authorization: Basic <api-token>
{
  "body": {
    "type": "doc",
    "version": 1,
    "content": [
      {"type": "paragraph", "content": [
        {"type": "text","text":"Quality v5"},
        {"type": "text","text":" Completeness 0.92, Clarity 0.81"}
      ]}
    ]
  }
}
\end{lstlisting}
\paragraph{GitHub check runs.} RFC repositories receive GitHub Check Runs summarizing completeness/clarity via \texttt{POST /repos/\{owner\}/\{repo\}/check-runs}:
\begin{lstlisting}[language=json]
{
  "name": "Requirement Quality",
  "head_sha": "abc123",
  "status": "completed",
  "conclusion": "neutral",
  "output": {
    "title": "Quality 0.84",
    "summary": "- Completeness: 0.90\n- Clarity: 0.78\n- Consistency: 0.85",
    "text": "Recommendations:\n1. Clarify acceptance criteria...\n2. Remove ambiguous terms..."
  }
}
\end{lstlisting}

\subsection{Slack Bots}
Slack Bolt bots notify reviewers when requirements move to “Review” or when quality drops. Buttons link back to the PWA, and slash commands fetch requirement snapshots. Webhooks originate from \texttt{/api/notifications/slack}.
\begin{lstlisting}[language=json]
POST https://slack.com/api/chat.postMessage
Authorization: Bearer xoxb-...
{
  "channel": "C012345",
  "blocks": [
    {"type":"section","text":{"type":"mrkdwn",
      "text":"*REQ-123 Ready for Review*\nQuality: 0.82"}},
    {"type":"actions","elements":[
      {"type":"button","text":{"type":"plain_text","text":"Approve"},
       "style":"primary","url":"https://pwa/requirements/REQ-123"},
      {"type":"button","text":{"type":"plain_text","text":"Request Changes"},
       "style":"danger","value":"request_changes"}
    ]}
  ]
}
\end{lstlisting}

\subsection{Experiment Platforms (Optimizely, LaunchDarkly)}
\begin{itemize}
    \item Optimizely experiments are created/updated through \texttt{POST /v2/experiments}. Webhooks feed results back to \texttt{/api/experiments/optimizely}.
\begin{lstlisting}[language=json]
POST https://api.optimizely.com/v2/experiments
Authorization: Bearer <optimizely-token>
{
  "project_id": 1234,
  "audience_ids": [6789],
  "key": "onboarding_activation",
  "primary_metric": "activation_rate",
  "variations": [
    {"key":"control"},
    {"key":"personalized_checklist"}
  ]
}
\end{lstlisting}
    \item LaunchDarkly flags are provisioned via \texttt{POST /api/v2/flags/\{proj\}/\{env\}} so rollout rules mirror requirement cohorts.
\begin{lstlisting}[language=json]
POST https://app.launchdarkly.com/api/v2/flags/default/onboarding-activation
Authorization: api-key
{
  "name": "Onboarding Activation",
  "key": "onboarding-activation",
  "variations": [
    {"value": false, "name": "control"},
    {"value": true, "name": "checklist"}
  ],
  "tags": ["experiment","req-123"]
}
\end{lstlisting}
\end{itemize}

\subsection{Automated Reporting to BI Tools}
\begin{itemize}
    \item Tableau: nightly jobs publish Hyper extracts using the Tableau REST/CLI APIs so executives can analyze quality + experiment metrics.
\begin{lstlisting}[language=bash]
python aggregate_metrics.py --output metrics.hyper
tableau-cli publish metrics.hyper "Bridging Worlds Metrics" \
  --server https://tableau.company.com --project "Product Insights"
\end{lstlisting}
    \item Power BI: streaming datasets ingest JSON rows via \texttt{POST /datasets/\{id\}/rows} to keep dashboards current.
\begin{lstlisting}[language=json]
POST https://api.powerbi.com/beta/datasets/{datasetId}/rows?key=...
[
  {
    "requirementId": "REQ-123",
    "qualityScore": 0.84,
    "metric": "activation_rate",
    "baseline": 0.18,
    "target": 0.21,
    "actual": 0.205,
    "updatedAt": "2025-11-20T12:00:00Z"
  }
]
\end{lstlisting}
\end{itemize}

By following this architecture, teams can deliver a PWA that embodies \textit{Bridging Worlds}: every chapter’s practices become interactive workflows, integrated metrics, and collaborative knowledge sharing inside a single installable experience.
